\section{The fourth-root rank bound and sequential approximation}

\begin{definition}[Uniform logarithmic coefficients]
    \label{def:aklv_accuracy_log_coefficients}
    \lean{TNLean.MPS.AreaLaw.aklvInitialRankLogCoefficient,
        TNLean.MPS.AreaLaw.aklvAccuracyLogCoefficient}
    \leanok
    For $A\ge1$, put
    \begin{align*}
        C_0(A)&=100(A+1)(8+2\log(A+1)),\notag\\
        C(A)&=10001C_0(A)+800000.
    \end{align*}
\end{definition}

\begin{lemma}[The initial scale]
    \label{thm:aklv_initial_log_rank_scale}
    \lean{TNLean.MPS.AreaLaw.aklv_initial_width_bounds}
    \leanok
    \uses{def:aklv_accuracy_log_coefficients,def:aklv_accuracy_rank_budget}
    Let $A\ge1$, $q\ge1$ and $0<g\le1$. Set
    $u=1+\log q+\log(1/g)$ and $s_0=\lceil Au^2/g\rceil$.
    Then $u,s_0\ge1$, $1/g,\log(1/g)\le s_0$,
    $T(g)\le1040s_0$, and
    \begin{align*}
        s_0&\le(A+1)u^2/g,&
        \log D_0&\le C_0(A)u^3/g.
    \end{align*}
\end{lemma}
\begin{proof}
    \leanok
    The ceiling is at least $Au^2/g$ and at most $Au^2/g+1$.
    Since $u\ge1$ and $g\le1$, the latter is at most $(A+1)u^2/g$.
    Apply the logarithm to the base of $D_0$, use $\log u\le u$,
    and multiply by its exponent $32(2s_0+1)$.
\end{proof}

\begin{lemma}[The two degree regimes]
    \label{thm:aklv_accuracy_two_degree_regimes}
    \lean{TNLean.MPS.AreaLaw.aklvAccuracyIncrement_lt_log_add_one,
        TNLean.MPS.AreaLaw.aklvAccuracyDegree_le_large_square,
        TNLean.MPS.AreaLaw.aklvAccuracyDegree_le_initial_square,
        TNLean.MPS.AreaLaw.aklvAccuracyIncrement_pos}
    \leanok
    \uses{def:aklv_accuracy_parameters,thm:aklv_initial_log_rank_scale}
    Suppose $0<g\le1$, $0<\rho\le1$, $1/g\le s_0$, and
    $T(g)\le1040s_0$. For the accuracy parameters $z,m,s,t,\ell$,
    one has $m\ge1$ and
    \begin{align*}
        z\le s_0&\implies \ell\le66s_0^2,\notag\\
        s_0\le z&\implies m<z+1\ \text{ and }\ \ell\le66m^2.
    \end{align*}
\end{lemma}
\begin{proof}
    \leanok
    Write $r=(z^3/g)^{1/4}$ and $W=2t+s$. In the first regime,
    $r\le s_0$, hence $m\le s_0$, and $W\le4225s_0$.
    Therefore $z\sqrt{W/g}\le65s_0^2$.
    In the second regime, $r\le z$, hence $m<z+1$, and
    $W\le4225z$. The identity $r^4=z^3/g$ gives
    $z\sqrt{W/g}\le65m^2$. Taking ceilings gives both bounds.
\end{proof}

\begin{lemma}[Logarithmic costs of a quadratic degree]
    \label{thm:aklv_quadratic_degree_log_rank}
    \lean{TNLean.MPS.AreaLaw.aklv_rank_exponent_le_of_degree_le_square,
        TNLean.MPS.AreaLaw.aklv_log_rank_le_initial_cost,
        TNLean.MPS.AreaLaw.aklv_log_rank_le_accuracy_cost}
    \leanok
    Suppose $q\ge2$, $1\le n\le s$, and $\ell\le66n^2$.
    Put $F=[q(\ell+1)]^{32(\lfloor\ell/s\rfloor+\lfloor\sqrt\ell\rfloor+1)}$.
    Then
    \begin{align*}
        \lfloor\ell/s\rfloor+\lfloor\sqrt\ell\rfloor+1&\le76n,\notag\\
        \log F&\le10000\log\bigl([q(n^2+1)]^{32(2n+1)}\bigr).
    \end{align*}
    If also $z\ge1$ and $n\le z+1$, then
    $\log F\le400000n(1+\log q+\log(z+1))$.
\end{lemma}
\begin{proof}
    \leanok
    The quotient is at most $66n$ and the square root at most $9n$.
    Bound $q(\ell+1)$ by $67q(n^2+1)$ for the first comparison,
    and by $67q(z+1)^2$ for the second. The inequality
    $\log q\ge\log2\ge1/2$ absorbs the numerical constants.
\end{proof}

\begin{lemma}[The uniform fourth-root logarithmic rank bound]
    \label{thm:aklv_accuracy_fourth_root_log_rank}
    \lean{TNLean.MPS.AreaLaw.aklvAccuracyRankBound_log_le_cost,
        TNLean.MPS.AreaLaw.aklvAccuracyRankBound_log_le_fourthRoot}
    \leanok
    \uses{thm:aklv_accuracy_rank_log,thm:aklv_accuracy_two_degree_regimes,
        thm:aklv_quadratic_degree_log_rank,thm:aklv_initial_log_rank_scale}
    In the accuracy-dependent argument of \cite[Section~7]{Arad2013AreaLaw},
    for $q\ge2$, $A\ge1$, $0<g\le1$ and $0<\rho\le1$, use
    $s_0=\lceil Au^2/g\rceil$ and the actual common rank budget $B$.
    Then
    \begin{align*}
        \log B&\le10001\log D_0+
            400000m(1+\log q+\log(z+1)),\notag\\
        \log B&\le C(A)\left[\frac{u^3}{g}+
            \left(\frac{z^3}{g}\right)^{1/4}
                (u+\log(z+1))\right].
    \end{align*}
\end{lemma}
\begin{proof}
    \leanok
    Split at $z=s_0$. Below that scale, absorb the filter cost into
    $10000\log D_0$. Above it, use the fourth-root increment in the
    second quadratic-degree estimate. The exact boundary rank costs at
    most $s_0\log q+2m\log q$. The first term is absorbed by
    $\log D_0$. Finally $r\ge1$ and $m<r+1\le2r$.
    The degree remains $\lceil z\sqrt{(2t+s)/g}\rceil$ throughout.
\end{proof}

\begin{lemma}[Sharp actual approximants at every physical cut]
    \label{thm:aklv_sharp_physical_per_cut}
    \lean{TNLean.MPS.AreaLaw.aklv_chain_dimension_and_gap_bounds,
        TNLean.MPS.AreaLaw.exists_aklv_sharp_physical_per_cut_approximation}
    \leanok
    \uses{thm:aklv_physical_accuracy_window,thm:aklv_accuracy_fourth_root_log_rank}
    There are universal $A\ge1$ and $C>0$ such that the following holds.
    Let $\Omega$ be a unit ground vector of a nearest-neighbour sum of
    supported positive contractions, at energy $\mu$, with positive
    rank-one projector gap $g$ and $\mu+g$ in its spectrum.
    Then $q\ge2$ and $g\le1$. For every $0<\rho\le1$ and every cut,
    the actual filter image or exact boundary alternative has squared
    error at most $\rho$, rank at most $B$, and the preceding logarithmic
    rank bound. No fitting window or upper-gap bound is assumed.
\end{lemma}
\begin{proof}
    \leanok
    The two distinct spectral values force $q\ge2$ and a nonempty chain.
    An endpoint unitary orthogonalizes the unique ground vector while
    changing the energy by at most one, giving $g\le1$.
    Apply the actual accuracy-window construction and its uniform rank
    comparison. The failed-window case retains the ground vector itself.
\end{proof}

\begin{lemma}[Discarded mass from a cut approximant]
    \label{thm:aklv_cut_approximation_tail}
    \lean{TNLean.MPS.AreaLaw.one_sub_kyFanNorm_le_of_unit_cut_rank_approximation}
    \leanok
    If $\Omega,y$ are unit chain vectors, $\|y-\Omega\|^2\le\rho$,
    and $\operatorname{SR}_k(y)\le D$, then the mass beyond the $D$
    leading eigenvalues of the left marginal of $\Omega$ is at most $\rho$.
\end{lemma}
\begin{proof}
    \leanok
    Expanding the squared distance gives
    $|\langle\Omega,y\rangle|^2\ge1-\rho$.
    The variational Schmidt-overlap inequality bounds that overlap by
    the sum of the $D$ largest marginal eigenvalues. If $D$ exceeds the
    marginal dimension, its discarded mass is zero.
\end{proof}

\begin{definition}[The actual sequential spectral approximant]
    \label{def:aklv_sequential_spectral_approximant}
    \lean{TNLean.MPS.AreaLaw.aklvSpectralTruncation,
        TNLean.MPS.AreaLaw.aklvNormalizedSpectralTruncation}
    \leanok
    Let $P_k$ be the projection onto the $D$ leading eigenvectors of
    the left marginal of $\Omega$. Apply their lifted projections in
    the order $L-1,L-2,\ldots,1$, obtaining $\Phi_0$.
    Put $\Phi=\Phi_0/\|\Phi_0\|$ when $\Phi_0\ne0$.
\end{definition}

\begin{theorem}[Sequential assembly and the sharp MPS bound]
    \label{thm:aklv_sharp_sequential_mps}
    \lean{TNLean.MPS.AreaLaw.aklvNormalizedSpectralTruncation_spec,
        TNLean.MPS.AreaLaw.exists_aklv_sharp_physical_sequential_mps_approximation}
    \leanok
    \uses{thm:aklv_sharp_physical_per_cut,thm:aklv_cut_approximation_tail,
        def:aklv_sequential_spectral_approximant,
        thm:chain_sequential_truncation_eigen,thm:chain_sequential_truncation_sq}
    For the preceding physical data and $0<\epsilon\le1$, choose
    \begin{align*}
        \rho&=\frac{\epsilon^2}{16(L+1)},&
        z&=\log\frac{2048(L+1)}{\epsilon^2},&D&=B.
    \end{align*}
    The actual vector $\Phi_0$ is nonzero, and its normalization is an
    open-boundary matrix product state of bond dimension at most $D$,
    with $\|\Phi-\Omega\|\le\epsilon$. Moreover
    \begin{align*}
        \log D\le C\left[\frac{u^3}{g}+
            \left(\frac{z^3}{g}\right)^{1/4}(u+\log(z+1))\right].
    \end{align*}
    The constants precede the chain, gap and accuracy. The actual
    excited spectral value is retained among the physical hypotheses.
\end{theorem}
\begin{proof}
    \leanok
    Each discarded mass is at most $\rho$. The general sequential
    projection inequality gives
    $\|\Phi_0-\Omega\|^2\le(L-1)\rho\le\epsilon^2/16$.
    Thus $\Phi_0\ne0$, and normalization at most doubles its distance
    to the unit vector $\Omega$. The same ordered product has every
    internal cut rank at most $D$ and hence an actual matrix product
    representation. Its scalar normalization preserves that bound.
\end{proof}
